\chapter{Digital Shadows versus ROS2 Systems: Verification and Analysis} \label{chap:sota}

Building on the fundamentals described in the previous chapter, the focus now shifts to an analysis of the preceding studies on creating Digital Shadows in robotics in \Cref{sec:DShadows}. \Cref{sec:createDShadows} introduces the conceptual model and workflow for Digital Shadows that will be utilised in the present thesis. Subsequently, the runtime verification of ROS2 systems, as it reflects another use case of dynamic system models, is examined in \Cref{sec:RVinROS2}, followed by an exploration of the analysis and monitoring of ROS2 systems in \Cref{sec:analyzingROS2}. Previous work on modelling ROS2 is described in \Cref{sec:modellingROS2}.


\section{Digital Twins and Digital Shadows in Previous Work} \label{sec:DShadows}
It is observed that the main focus of Digital Shadows and Digital Twins is on the industry \cite[p.~375]{digitaltwinmodelling} and has a growing interest in robotics \cite[p.~4]{mazumder2023towards}. This section presents several approaches that arise from different robotic applications.

\Textcite{IndustrialRobotDigitalTwin} proposes a framework for the development of an Industrial Robot Digital Twin. It develops a comprehensive ontology model, encompassing structural and functional data such as state information and measurement values. Furthermore, the approach encompasses static information, such as names and types, as well as dynamic, time-dependent information, including movement, position, and load.
To obtain an accurate Digital Twin instance, it covers the entire lifecycle of the industrial robot, reaching from the design phase to manufacturing, utilisation, and maintenance to recycling. Consequently, the information from each phase influences the Digital Twin instance, which combines model-based and simulation-based techniques. This facilitates real-time monitoring and predictive maintenance --~a proactive maintenance concept that utilises continuous monitoring and data analysis to predict the condition of a system. The model is based on \ac{mbse} using \ac{sysml}, ensuring that all aspects of the system are modelled consistently and comprehensively. In order to enable different variants, the paper utilises a feature model that categorises features into mandatory, optional, and alternative features, thus enabling flexibility while maintaining a basic structure.

\Textcite{DSforLowDelayAnalysisOfProdQualityRisk} proposes a Digital Shadow to support \acf{fmea} --~a systematic method for identifying and evaluating potential error sources in a system~-- in robotic systems on the shop floor. \ac{fmea} experts and developers collaborate to link \ac{fmea} models with production, process, and resource data, forming a comprehensive system that maps various types of faults and their effects on the running system. This knowledge is structured and stored inside a knowledge graph, which represents the virtual model and refers to the physical system. Afterwards, experts select relevant data sources and variables to capture the physical system's characteristics and continuously update the virtual model. This facilitates the monitoring of the running system's performance over a specified time period.

\Textcite{ROSIE} proposes ROSIE, a \acs{ros} adapter that aims to provide a plug-and-play solution to create a Digital Twin for ROS1 systems. The primary objective of this proposal is to facilitate the operation and configuration of industrial robot arm movements by providing two types of goals that can be set using augmented reality. The first goal is to move the robot's arm to an arbitrary position in space. The second is to set up a trajectory that the robot's arm shall follow. The integration is achieved by augmenting the \acs{ros}-control package with an additional abstraction layer. The \acs{ros}-control package provides access to robotic actuators and consists of several abstraction layers to fit agnostic, \acs{ros}-based robots. By adding another abstraction layer, \Textcite{ROSIE} achieves robot-independent control of the physical system. The virtual model is updated by subscribing to the \glqq{}Joint State\grqq{} topic, which contains the latest information about the joint states. It is important to note that ROSIE is designed for ROS1. However, the authors are currently working on migrating the tool to ROS2.

In summary, \Textcite{IndustrialRobotDigitalTwin}, \Textcite{DSforLowDelayAnalysisOfProdQualityRisk}, and \Textcite{ROSIE} offer valuable contributions, yet each is tailored to a specific industrial context. \Textcite{IndustrialRobotDigitalTwin} presents an Industrial Robot Digital Twin that spans the entire lifecycle and utilises a feature-driven approach. However, the need to predefine features would limit flexibility for broader ROS2 applications and thus contradicts to the objectives of this thesis. Another focus is on predictive maintenance and therefore the aim is to prevent failures rather than to detect and diagnose faults. In contrast, \Textcite{DSforLowDelayAnalysisOfProdQualityRisk} focuses on the production phase, integrating \acl{fmea} to evaluate product quality. \Textcite{ROSIE} employs a distinct approach by leveraging the ROS-control package to mirror physical systems into augmented reality, facilitating the deployment of industrial robots. However, it restricts flexibility and does not specifically address \acl{fdd} needs, thereby contradicting the objectives of this thesis. \newline
Multiple elements, such as the division of static and dynamic data in \Textcite{IndustrialRobotDigitalTwin} and \Textcite{DSforLowDelayAnalysisOfProdQualityRisk}, the utilisation of SysML for modelling in \Textcite{IndustrialRobotDigitalTwin}, and the utilisation of a knowledge graph in \Textcite{DSforLowDelayAnalysisOfProdQualityRisk} are promising. Nevertheless, these approaches do not fully satisfy the requirement for a flexible, ROS2-wide solution that facilitates \acl{fdd} in the development phase during rapid prototyping.

To create a Digital Shadow that matches our requirements, the subsequent section presents a conceptual model for Digital Shadows and a workflow for creating a Shadow Instance.

\section{Creation of a Digital Shadow} \label{sec:createDShadows}
In accordance with the descriptions provided in \Cref{sec:DMDSDT} and the required automated data flow from the physical object to the virtual model exclusively, the Digital Shadow proves to be a suitable solution that meets our requirements. Furthermore, this thesis adopts the following definition:
\glqq{}A digital shadow is a set of temporal data traces and/or their aggregation and abstraction collected concerning a system for a specific purpose with respect to the original system.\grqq{} \cite[p.~3]{DSDTDefinition}.


    \subsection{Conceptual Model of a Digital Shadow} \label{subsec:DSConceptualModel}
    The conceptual model utilised in this thesis is informed by the conceptual model proposed by \Textcite{DSConceptualModelElsevier} and \Textcite{DSConceptualModelSpringer}, and illustrated in \Cref{fig:tikz:DSconceptualModel}. In general, it specifies which concepts and relations a structure requires to function as a Digital Shadow.

    \begin{figure}[h]
        \centering
        % \includegraphics{figures/DSConceptual.png}
        \begin{tikzpicture}[>={Triangle[angle=60:4]},help help lines/.style={lightgray!50,very thin}, help lines/.style={gray!50,thin}]
    \node[umlClass] at (0,0)    (shadow)    {Digital Shadow};
    \node[umlClass] at (0,-1.5) (tracer)    {DataTracer};

    \node[umlClass] at (4, .5) (purpose)    {Purpose};
    \node[umlClass] at (4,-.5) (model)      {Model};

    \node[umlClass] at (-4,.5) (asset)      {Asset};

    \node[umlClass] at (-4,-1) (source)     {Source};

    \node[umlClass] at (-7,-1) (property)   {Property};
    \node[umlClass] at (-7,-2) (measurement){Measurement};
    \node[umlClass] at (-7,-3) (processing) {Processing};
    \node at (-7, -2.5) {...};

    \node[umlClass] at (-1.25,-3) (point)   {DataPoint};
    \node[umlClass] at ( 1.25,-3) (metadata){MetaData};

    \node[umlClass] at (2.8,-1.75) (structure) {Structure};
    \node[umlClass] at (5.2,-1.75) (behavior)  {Behavior};
    \node at (4, -1.9) {...};

    \node[umlClass] at (4,-3.25) (sysconf)   {System\\Configuration};

    %%% Compositions
    \draw [-diamond] (tracer.north) -- node [at start, above right = -.1cm and 0cm] {\footnotesize 1..*} (shadow.south);
    \draw [-diamond] (property.east) -- node [at start, above right] {\footnotesize *} (source.west);
    \draw [diamond-, rounded corners] (tracer.south) ++(-.5,0) -- ++(0,-.5) -| node [at end, above right = -.1cm and 0cm] {\footnotesize 1..*} (point.north);
    \draw [diamond-, rounded corners] (tracer.south) ++( .5,0) -- ++(0,-.5) -| node [at end, above left = -.1cm and 0cm] {\footnotesize 1..*} (metadata.north);
    \draw [-diamond, rounded corners] (asset.north) |- node [at start, above right] {\footnotesize *} ++ (-1.1,.25) |- (asset.west);

    %%% Inheritances
    \draw [-open triangle 60] (asset.south) -- (source.north);
    \draw [-open triangle 60, rounded corners] (measurement.east) -| (source.south);
    \draw [-open triangle 60, rounded corners] (processing.east) -| (source.south);
    \draw [-open triangle 60, rounded corners] (structure.north) -- ++(0,.285) -| (model.south);
    \draw [-open triangle 60, rounded corners] (behavior.north) -- ++(0,.285) -| (model.south);
    \draw [-open triangle 60, rounded corners] (sysconf.north) |- ++(-2.4,.3) -- ++(0,1.205) -| (model.south);


    %%% Association
    \draw [->, rounded corners] (shadow.north) ++(-1,0) -- ++(0,.25) -- node [midway, above] {\footnotesize links to} ++(2,0) --  node [at end, above right] {\footnotesize *} ++(0,-.25);
    \draw [<-, rounded corners] (sysconf.west) ++(0,-.25) -| node [midway, below right = -.1cm and -.25cm] {\footnotesize knows} node [at start, below left] {\footnotesize 1} (metadata.south);
    \draw [->, rounded corners] (shadow.west) -- ++(-1,0) |- node [midway, above right = 0cm and -.5cm] {\footnotesize stands for} node[at end, above right] {\footnotesize 1} (asset.east);
    \draw [->, rounded corners] (tracer.west) -- ++(-1,0) |- node [midway, above right = -.05cm and -1cm] {\footnotesize originates from} node [at end, below right] {\footnotesize 1} (source.east);
    \draw [->, rounded corners] (shadow.east) ++(0, .15) -- ++(1,0) node [midway, above] {\footnotesize fulfills} |- node [at end, below left] {\footnotesize 1} (purpose.west);
    \draw [->, rounded corners] (shadow.east) ++(0,-.15) -- ++(1,0) node [midway, below] {\footnotesize uses} |- node [at end, above left] {\footnotesize 1..*} (model.west);

\end{tikzpicture}
        \caption[Conceptual Model of a Digital Shadow]{Conceptual Model of a Digital Shadow according to \cite{DSConceptualModelElsevier}\cite{DSConceptualModelSpringer}} \label{fig:tikz:DSconceptualModel}
    \end{figure}

    According to the established definition and as illustrated in \Cref{fig:tikz:DSconceptualModel}, each Shadow Instance is characterised by a distinct Purpose. The \emph{Purpose} serves to filter the data, ensuring efficiency by maintaining focus on relevant information.
    The \emph{Model} is characterised by, e.g., the behaviour and structure of the Asset and is utilised as a construction plan. Guided by the Purpose, the Model should only represent the relevant aspects of the system.
    The \emph{Asset} reflects the object that is to be mirrored by the Shadow Instance. It may include subsystems that are classified as Assets. It is important to mention that each Shadow Instance is linked to a single Asset, while an Asset can be reflected by multiple Shadow Instances.
    Establishing suitable attributes that accurately represent the properties of the Asset facilitates their utilisation as a data \emph{Source} and therefore the creation of a Shadow Instance.
    Each property or set of properties represented by the Shadow Instance is obtained by tracing a specific Source utilising a connected \emph{DataTracer}.
    A \emph{DataPoint} represents a single value from a single DataTracer. Since each DataTrace can be associated with several DataPoints, they can, e.g., reflect a time series.
    To contextualise DataTraces, they hold \emph{MetaData} that describes them and holds information about the \emph{System Configuration}. \cite[p.~273ff]{DSConceptualModelSpringer}

    Based on the conceptual model, the following section focuses on the workflow of creating a Shadow Instance.

    \subsection{Workflow for Creating a Shadow Instance} \label{subsec:DSWorkFlow}

    \Cref{fig:tikz:DSworkflow} illustrates the process of generating a Shadow Instance based on multiple data inputs and predefined structures.
    \begin{figure}[h]
        \centering
        \begin{tikzpicture}[>={Triangle[angle=60:4]},help help lines/.style={lightgray!50,very thin}, help lines/.style={gray!50,thin}]

    \tikzset{
        ExecutingBlock/.style={
            rectangle,
            rounded corners = 5,
            draw = black,
            inner sep = 5,
        },
        Block/.style={
            rectangle,
            draw = black,
            inner sep = 5,
        },
    }

    %%% Nodes
    \node [Block] (instance) {Shadow \textbf{Instance}};
    \node [ExecutingBlock, above = .5cm of instance] (caster) {Shadow \textbf{Caster}};
    \node [above = 1.4cm of caster] (data) {\textbf{Data}};
    \node [Block, below left = .5cm of data] (type) {Shadow \textbf{Type}};
    \node [Block, below right = .4cm and .9cm of data] (models) {\textbf{Models}};

    \node [above left = .5cm and .9cm of data] (asset) {\textbf{Asset}};
    \node [above right = .4cm and 0cm of data] (others) {\textbf{other Systems}}; % , text width = 1.75cm

    %%% Arrows
    \draw [->] (caster) -- (instance);
    
    \draw [->] (type.south) |- (caster.west);
    \draw [->] (data.south) -- (caster.north);
    \draw [->] (models.south) |- (caster.east);

    
    \draw [->] (asset.south) |- (data.west);
    \draw [->] (others.south) |- (data.east);

\end{tikzpicture}
        \caption[Workflow of a Digital Shadow]{Workflow of a Digital Shadow according to \cite{DSConceptualModelElsevier}} \label{fig:tikz:DSworkflow}
    \end{figure}

    The \emph{Asset}, as described in \Cref{subsec:DSConceptualModel}, represents the specific object for which a Shadow Instance is created. Additionally, \emph{other Systems}, such as external simulations or data sources, contribute data that, while not directly linked to the Asset, may still influence its behaviour.
    Both the Asset and other Systems transmit data to the \emph{Shadow Caster}, which is responsible for constructing the Shadow Instance. This process is guided by two key elements: the Shadow Type and further Models. The \emph{Shadow Type} functions as a blueprint, defining the general structure of the Shadow Instance based on its intended purpose and ensuring consistency across similar entities. In contrast, \emph{Models}, formulating domain knowledge, provide object-specific details essential for refining the generated Shadow Instance.\newline
    Together, these components enable the Shadow Caster to create an accurate Shadow Instance, tailored to the characteristics and environmental influences affecting the Asset. \cite[p.~836]{DSConceptualModelElsevier}

    The presented conceptual model has been employed successfully in a number of works, including \Textcite{IIOTinsightsByDS} and \Textcite{DSinProdControl}.

This section provided a comprehensive overview of the creation of a Shadow Instance and the model of a Digital Shadow. As this thesis focuses on supporting \acl{fdd} on ROS2 systems and thus facilitating the analysis of such systems, the subsequent sections explore multiple approaches to validate and analyse ROS2 systems.


\section{Runtime Verification in ROS2 Systems} \label{sec:RVinROS2}
Runtime verification is a verification technique aimed at checking the correctness of a system during its execution. The core idea of runtime verification is to instrument a system with monitors that observe execution traces and detect violations or deviations from the expected behaviour. Therefore, runtime verification checks whether the functionality of a system meets the requirements and should not be confused with runtime validation, which focuses on whether those requirements are compatible with the system's intended purpose.

\Textcite{ROSRV} introduces ROSRV, a runtime verification framework for \ac{ros}-based robotic systems, aiming to enhance safety and security by monitoring and controlling system interactions in real time. Through a specification language, users can define safety properties and security policies, which are then translated into automatically generated monitors that observe message flows and trigger corrective actions if necessary. To achieve this, ROSRV, designed for ROS1, replaces the ROS Master with a so-called RVMaster.

\Textcite{ROSMonitoring} and \Textcite{ROSMonitoring2.0} introduce ROSMonitoring and ROSMonitoring 2.0, respectively. Both share the common goal of enhancing runtime verification in \ac{ros}-based systems, ensuring that robotic applications adhere to predefined constraints during execution. They operate by monitoring system behaviour during runtime, rely on formal specifications to define correct behaviour, and aim to improve system reliability and safety. Additionally, both frameworks integrate automated monitoring mechanisms, reducing manual effort in verifying system compliance. Neither approach alters the core functionality of \ac{ros} but instead adds a verification layer that passively or actively ensures constraints are met. Furthermore, both systems have been evaluated in real-world robotic scenarios to analyse their scalability and overhead impact.

Despite these similarities, the two approaches differ in implementation and scope. ROSMonitoring 2.0 focuses on automating the generation of runtime monitors from natural language requirements, enabling developers to specify constraints in a structured format that is automatically transformed into executable monitoring components. These monitors passively observe system behaviour, identifying deviations from expected conditions but without direct intervention. In contrast, ROSMonitoring actively validates message exchanges, effectively enforcing constraints rather than merely detecting violations.
Another distinction lies in flexibility and integration. While ROSMonitoring 2.0 follows a structured workflow with a dedicated pipeline for transforming requirements into monitors, ROSMonitoring enables users to define verification properties using diverse specification formalisms, making it more adaptable for differing use cases.
Despite their strengths in constraint validation, neither approach provides a comprehensive \acl{fdd} mechanism, as they focus on verifying rule compliance rather than analysing failure causes. While both enhance system robustness by ensuring correct operation within defined constraints, a more holistic model-driven approach would be required to systematically detect, diagnose, and mitigate faults in ROS2-based robotic systems.

In summary, ROSRV \cite{ROSRV}, ROSMonitoring \cite{ROSMonitoring} and ROSMonitoring 2.0 \cite{ROSMonitoring2.0} focus on runtime verification and thus prevent hazardous behaviour based on rules. The definition of these rules is achieved through the utilisation of differing languages, necessitating deep expert knowledge. In contrast, this thesis introduces a meta-model that reduces the need for expert knowledge, facilitates \acl{fdd}, and can be used with various algorithms, including those that perform runtime verification.


\section{Analysis of ROS2 Systems} \label{sec:analyzingROS2}
Given the present thesis emphasis on developing a meta-model that supports \acl{fdd} by facilitating the analysis of the ROS2 system, the focus now shifts to examining prior research conducted on analysing ROS2 systems. Accordingly, this section also explores several implementation specific details that will be particularly relevant for the evaluation of the approach proposed in this thesis (see \Cref{chap:evaluation}).
The research leads to several works that focus on analysing ROS2 system performance, particularly in terms of latency measurements \cite{howFastIsMySoftware}\cite{ROS2exploringPerformance}\cite{multiExecutor}\cite{endToendLatency}.

\emph{ros2\_tracing} \cite{ros2_tracing} is a low-overhead tracing tool that employs several tracepoints within the ROS2 source code. These tracepoints provide detailed information about events occurring within a ROS2 system, such as the initialisation of nodes and publishers \cite{github:ros2_tracing}. 
At the time of this work, ros2\_tracing only enables offline analyses and is therefore unsuitable for the objectives of this thesis, as the monitoring should occur at runtime. However, the tracepoints are included within the ROS2 source code and enabled to be traced, and thus analysed, online. Independent of ros2\_tracing.\newline
Multiple previous works demonstrate the potential of ros2\_tracing. For instance, \Textcite{howFastIsMySoftware} and \Textcite{multiExecutor} employed ros2\_tracing to collect fine-grained runtime data and identify critical timing bottlenecks, while evaluating the impact of different system configurations, such as varying executor strategies. \Textcite{ros2_tracing:msgFlowAnalysis} utilises ros2\_tracing to analyse the message flow throughout a ROS2 system.\newline
For the objectives of this thesis, it is imperative to mention that tracepoints only generate a latency overhead when actively traced. The impact of tracepoints related to publishing messages with differing sizes (1~KiB, 32~KiB, 64~KiB, 256~KiB) and frequencies (100~Hz, 500~Hz, 1000~Hz, 2000~Hz) was measured by \Textcite[p.~6516f]{ros2_tracing}. The mean latency overhead quantifies to 3.3~$\upmu$s and therefore the created overhead \glqq{}is acceptable for most target situations in ROS2.\grqq{} \cite[p.~6516f]{ros2_tracing}. Based on this evaluation, and also due to the high number of tracepoints, tracing those tracepoints within the ROS2 source code offers a promising approach for analysing ROS2 systems.

\Textcite{ROS2exploringPerformance} employed iRobot's \emph{ros2-performance} \cite{github:ros2_performance} to measure and analyse latencies. Since the framework is primarily designed to evaluate single-process applications offline, it does not fit our requirements.

Another approach for analysing a ROS2 network is the utilisation of the \emph{ros2-diagnostics} framework. The framework's functionality is restricted to the provision of tools for standardised messages and evaluation. However, the establishment of the integration interface on the Node side is a responsibility that falls to the user. Consequently, it does not fully meet the objectives of this thesis, as this constitutes an intervention into the codebase \cite{github:ros_diagnostics}.

In summary, the most promising contribution is provided by ros2\_tracing \cite{ros2_tracing}, as it enables a non-intrusive method to observe ROS2 systems by tracing specific tracepoints. Alternatives for online analysis, such as ros2-diagnostics \cite{github:ros_diagnostics}, are intrusive and therefore necessitate modification of the codebase to observe a ROS2 system in question.

Although runtime verification and system analysis provide valuable insights into the behaviour and correctness of ROS2 systems during execution, these approaches either rely on deep expert knowledge or analyse only fractions of ROS2 systems. However, effective \acl{fdd} in ROS2 systems necessitates an approach that considers the system as a whole. ROS2 system models offer such a comprehensive perspective.


\section{Modelling ROS2 Systems} \label{sec:modellingROS2}
Recent research has demonstrated the significant role of models in leveraging Digital Twins and Digital Shadows. Nevertheless, as the objective is to create a Digital Shadow that facilitates \acl{fdd} for generic ROS2 systems, this thesis aims to create a meta-model depicting ROS2 systems in general. In the literature, such models are seldom encountered in the context of ROS2.

One approach is presented by \emph{MeROS} \cite{MeROS}, a \ac{sysml}-based meta-model developed for both \ac{ros} versions. MeROS systematically encapsulates core \ac{ros} concepts --~including Nodes, Topics, Services, and Actions~-- by grouping related communication methods and components into abstract entities. This approach not only improves the clarity and compactness of system descriptions but also facilitates system documentation, defect detection, and team collaboration.
Its design is driven by a comprehensive set of requirements derived from both literature and practical experiences in ROS-based projects. However, while MeROS excels in providing a unified model for the architecture and behaviour of ROS systems, it primarily targets the areas of system documentation and integration. Therefore, it provides guidance for \ac{ros} systems that are in the preliminary stages of their design process.
Notably, it does not offer explicit constructs or extensions aimed at \acl{fdd}. Furthermore, the created model refers to ROS2 only and does not consider the underlying operating system, which may also be prone to faults.

Further approaches are designated for ROS1 only and lack completeness, as concepts such as Actions might not be represented.
For instance, \emph{ROSMOD} \cite{ROSMOD} is a model-driven development environment for distributed \ac{ros} systems, focusing on modelling, code generation, and deployment. While it supports runtime monitoring, it lacks direct \acl{fdd} capabilities. In contrast, \emph{RoBMEX} \cite{RobMex} is a domain-specific modelling framework for graphical mission planning and automatic code generation in \ac{ros}-based systems. It does not integrate runtime monitoring and is designed for mission specification rather than system diagnosis.

To date, the literature lacks a dedicated meta-model for ROS2 systems that is tailored to \acl{fdd}, leaving a notable gap in the state of the art.
This gap motivates the development of a meta-model for ROS2 systems that integrates mechanisms which support \acl{fdd}, paving the way for more resilient and reliable rapid prototypes. Therefore, the subsequent chapter employs \acl{mbse} to model various concepts and aspects of ROS2 systems.